\chapter{Основная часть}

\section{Формализованная постановка задачи}

Во время выполнения выпускной квалификационной работы была формализована  в виде диаграммы IDEF0 постановка задачи коммивояжёра. Формализованная постановка задачи представлена на рисунке~\ref{tsp-idef0}.

\begin{figure}[H]
	\centering
	\includegraphics[width=\textwidth]{tsp/01\_A0}
	\caption{Формализация задачи коммивояжёра в виде IDEF0-диаграммы}
	\label{tsp-idef0}
\end{figure}

На рисунках~\ref{fig:diplom}-\ref{fig:03A4} представлена формализация разработанного метода в виде детализированной диаграммы IDEF0.

\begin{figure}[H]
	\centering
	\includegraphics[width=\textwidth]{method/01\_A0}
	\caption{Формализация предлагаемого метода в виде IDEF0-диаграммы}
	\label{fig:diplom}
\end{figure}

\begin{figure}[H]
	\centering
	\includegraphics[width=\textwidth]{method/02\_A0}
	\caption{IDEF0 диаграмма метода 1-го уровня}
	\label{fig:02A0}
\end{figure}

\begin{figure}[H]
	\centering
	\includegraphics[width=\textwidth]{method/03\_A2}
	\caption{IDEF0 диаграмма формирования путей муравьёв}
	\label{fig:03A2}
\end{figure}

\begin{figure}[H]
	\centering
	\includegraphics[width=\textwidth]{method/03\_A4}
	\caption{IDEF0 диаграмма формирования новой популяции}
	\label{fig:03A4}
\end{figure}


\section{Функциональное тестирование}


На рисунке~\ref{fig:test} представлена схема алгоритма функционального тестирования для одного теста.

\begin{figure}[H]
	\centering
	\includegraphics[width=0.6\textwidth]{test}
	\caption{Схема алгоритма функционального тестирования}
	\label{fig:test}
\end{figure}

Для тестирования были выделены следующие классы эквивалентности входных графов и параметров:
\begin{itemize}
	\item граф малой размерности ($n < 10$ );
	\item граф средней размерности ($10 \leqslant  n < 100 $);
	\item граф большой размерности ($n \geqslant 100$);
	\item граф заданный в двумерной евклидовой форме (EUC\_2D);
	\item граф заданный в трёхмерной евклидовой форме (EUC\_3D);
	\item граф заданный географическими координатами (GEO);
	\item граф заданный явной матрицей весов (EXPLICIT);
	\item все возможные тройки стратегий селекции, скрещивания, мутации.
\end{itemize}

Функциональное тестирование проведено для всех выделенных классов эквивалентности. Во всех тестовых запусках:
\begin{itemize}
	\item сформированный маршрут являлся гамильтоновым циклом;
	\item алгоритм завершал работу без ошибок;
	\item длина маршрута соответствовала ожидаемому диапазону значений;
	\item параметры $\alpha$ и $\beta$ менялись по ходу маршрута индивидуально у каждого муравья. 
\end{itemize}

Все функциональные тесты успешно пройдены.

\section{Модульное тестирование}

Модульное тестирование приложения проводилось с помощью стандартного пакета Go -- testing, и стороннего пакета testify.

В качестве меры качества покрытия кода использовалась метрика стандартной утилиты cover входящей в состав Go, которая измеряет покрытие кода как процент от выполненных хотя бы единожды выражений.

Тестами были покрыты модули стратегий, структур данных, а также части ядра, посвященные вычислениям.

Покрытие кода составило \textbf{42.2\%}.

\section{Тестирование графического интерфейса}

Тестирование графического интерфейса не проводилось.

\section{Исследование характеристик разработанного программного обеспечения}

\subsection{Исследование эволюционных операторов}

\textbf{Постановка и описание}

Объектом данного исследования является разработанный муравьиный алгоритм с динамическими гетерогенными параметрами на основе генетического алгоритма. Предметом исследования является зависимость качества получаемых алгоритмом решений от конфигурации эволюционных операторов. Цель исследования – определить оптимальную конфигурацию эволюционных операторов, обеспечивающую наилучшее качество решений задач коммивояжёра.

Варьируемыми факторами исследования являются: стратегия селекции (элитарная, турнирная, рулеточная), стратегия скрещивания (арифметическая, BLX, SBX), стратегия мутации (равномерная, гауссовская). Также варьируется входной взвешенный неорграф по 6 задачам в формате TSPLIB. 

Постоянными параметрами являются: начальные значения $\alpha=1$ и $\beta=3$, коэффициент испарения $\rho=0.5$, множитель феромона $Q=1$, количество муравьёв $m=n$, начальный уровень феромона $\tau_0 = m / C_n^m$, количество поколений 300, период адаптации $w=5$, количество родителей $0.4m$. Для сравнения конфигураций используются средняя относительная ошибка и средний ранг.

\textbf{Результаты}

На рисунке~\ref{fig:variants-bar} представлена столбчатая диаграмма средних рангов конфигураций, а на рисунке~\ref{fig:variants-heat} -- тепловая карта рангов конфигураций по входным графам.

\begin{figure}[H]
	\centering
	\includegraphics[width=\textwidth]{variants/rank}
	\caption{Диаграмма средних рангов конфигураций эволюционных операторов}
	\label{fig:variants-bar}
\end{figure}

\begin{figure}[H]
	\centering
	\includegraphics[width=0.8\textwidth]{variants/heatmap}
	\caption{Тепловая карта рангов конфигураций эволюционных операторов по входным графам}
	\label{fig:variants-heat}
\end{figure}

По среднему рангу и средней относительной ошибке наиболее удачной конфигурацией является сочетание рулеточной селекции, BLX-скрещивания и гауссовской мутации. Данный вариант использован в дальнейших исследованиях как основной.

\subsection{Исследование сходимости}

\textbf{Постановка и описание}

Объектом данного исследования является разработанный гетерогенный муравьиный алгоритм с оптимальными эволюционными операторами. Предметом исследования является сходимость алгоритма к субоптимальным решениям в зависимости от количества поколений и наличия локальной оптимизации. Цель исследования – оценить влияние количества поколений и локальной оптимизации на сходимость разработанного алгоритма, а также сравнить с классическим муравьиным алгоритмом.

Варыруемыми факторами являются: количество поколений (500, 800, 1000) и тип алгоритма (классический муравьиный алгоритм, разработанный алгоритм без локальной оптимизации, разработанный алгоритм с 2-opt локальной оптимизацией). Исследование проводилось на 7 задачах TSPLIB.

Постоянные параметры: $\alpha=1$, $\beta=5$, $\rho=0.5$, $Q=1$, $m=n$, $\tau_0 = m / C_n^m$, $w=5$, количество родителей $0.4m$. Для оценки сходимости используются количество итераций до лучшего решения (ITB), длина наилучшего маршрута и затраченное время.

\textbf{Результаты}

На рисунке~\ref{fig:convergence-boxplot} представлена коробчатая диаграмма, отображающая распределение медианы ITB полученным алгоритмами в задачах, а на рисунке~\ref{fig:convergence-scatter} -- точечный график зависимости медиан длины маршрута и ITB разными алгоритмами по задачам.

\begin{figure}[H]
	\centering
	\includegraphics[width=\textwidth]{convergence/boxplots}
	\caption{Коробчатая диаграмма распределения ITB по задачам и алгоритмам}
	\label{fig:convergence-boxplot}
\end{figure}

\begin{figure}[H]
	\centering
	\includegraphics[width=\textwidth]{convergence/scatter}
	\caption{Точечный график зависимости ITB и длины найденного маршрута по задачам и алгоритмам}
	\label{fig:convergence-scatter}
\end{figure}

Разработанный алгоритм без локальной оптимизации демонстрирует качество, сопоставимое с классическим (разница 1–2\%), но достигает его в 3.7–8 раз быстрее по числу итераций. Добавление 2-opt локальной оптимизации улучшает качество решений на 6–8\%, однако увеличивает время выполнения в 1.3–1.7 раза.

\subsection{Исследование устойчивости}

\textbf{Постановка и описание}

Объектом данного исследования является разработанный гетерогенный муравьиный алгоритм с оптимальными эволюционными операторами. Предметом исследования является зависимость качества решений от начальных значений параметров $\alpha$ и $\beta$. Цель исследования – оценить устойчивость разработанного алгоритма к выбору параметров $\alpha$ и $\beta$ в сравнении с классическим муравьиным алгоритмом.

Варыруемыми факторами являются: начальное значение $\alpha$ от 0.2 до 2.0 с шагом 0.2, начальное значение $\beta$ от 0.5 до 5.0 с шагом 0.5, тип алгоритма (классический и разработанный).

Постоянные параметры: $\rho=0.5$, $Q=1$, $m=n$, $\tau_0 = m / C_n^m$, количество поколений 500, $w=5$, количество родителей $0.4m$. Для оценки устойчивости используются размах средних длин маршрута и коэффициент вариации.

\textbf{Результаты}

На рисунке~\ref{fig:stability-diff} изображены графики зависимости средних длин маршрутов полученных муравьиным алгоритмом и разработанной модификацией от параметров $\alpha$ и $\beta$.

\begin{figure}[H]
	\centering
	\includegraphics[width=\textwidth]{stability/comparison}
	\caption{Сравнение средних длин маршрутов, полученных муравьиным алгоритмом и разработанной модификацией при варьировании $\alpha$ и $\beta$}
	\label{fig:stability-diff}
\end{figure}

Разработанная модификация демонстрирует значительно более высокую устойчивость: коэффициент вариации средней длины маршрута снижен с 27,96\% до 4,54\%, а размах – со 1101,85 до 124,95 единиц. Предложенный алгоритм сохраняет высокое качество решений независимо от параметров, в то время как классический требует более тонкой настройки.

\section*{Вывод}